\documentclass{article}

% NeurIPS 2026 style, preprint option: no line numbers, authors shown.
\usepackage[preprint]{neurips_2026}

\usepackage[utf8]{inputenc}
\usepackage[T1]{fontenc}
\usepackage{hyperref}
\usepackage{url}
\usepackage{booktabs}
\usepackage{amsfonts}
\usepackage{nicefrac}
\usepackage{microtype}
\usepackage{xcolor}
\usepackage{graphicx}
\usepackage{adjustbox}
\usepackage{array}
\usepackage{enumitem}
\usepackage[most]{tcolorbox}
\usepackage[capitalise,noabbrev]{cleveref}

\input{generated/tables/macros}

% Question box at the top of each analysis section.
\newtcolorbox{question}{enhanced, breakable, colback=white, colframe=black!35,
  boxrule=0.4pt, leftrule=2pt, arc=0pt, left=6pt, right=6pt, top=3pt, bottom=3pt,
  fontupper=\small}
% Analysis still to write (shown in color so that it can't be mistaken for results).
\newcommand{\placeholder}[1]{{\color{orange!80!black}\textit{[To write: #1]}}}
% Figure that needs data not collected yet (all of it, or part of it, synthetic).
% Shown only when the figure was drawn with synthetic data (flags written by analyze.py plot).
\IfFileExists{generated/figures/synthetic.tex}{\input{generated/figures/synthetic.tex}}{}
\makeatletter
\newcommand{\syntheticnote}[1]{\expandafter\ifx\csname syn@#1\endcsname\relax\else
  \if1\csname syn@#1\endcsname{\color{red!70!black}\textbf{Synthetic data}: placeholder values from a
  toy model, not measurements; hollow markers.}\ \fi\fi}
\makeatother
% Generated figure, or a box saying its data is missing (the figure file is deleted when it has no data).
\newcommand{\fig}[2][\linewidth]{\IfFileExists{generated/figures/#2.pdf}%
  {\includegraphics[width=#1]{generated/figures/#2}}%
  {\fbox{\parbox{0.9\linewidth}{\centering\small\vspace{1em}No data yet for \texttt{\detokenize{#2}}.\vspace{1em}}}}}

\title{Large-scale LLM pre-training on Rorqual:\\ hardware, fabric, scaling and readiness}

\author{%
  Milabench team \\
  Mila -- Quebec AI Institute \\
}

\begin{document}

\maketitle

\begin{abstract}
  We evaluate Rorqual (Alliance) for large-scale LLM pre-training, as a test bed for the evaluation
  workflow of upcoming exascale clusters. Two workloads are used: Qwen3.8-27B (dense, TorchTitan,
  FSDP over all GPUs) and GLM-5 (mixture of experts, Megatron). We report node health, the cost of the
  network and of node placement on its leaf switches, weak scaling up to the full machine, MoE-specific
  behaviour, software-stack efficiency, and operational readiness.
  \placeholder{key numbers once the October~30 window is analyzed.}
\end{abstract}

\begin{tcolorbox}[colback=red!4, colframe=red!50!black, boxrule=0.4pt, arc=0pt, fontupper=\small]
  \textbf{Draft generated \DataGenerated.} \NRunsReal{} runs from the preliminary campaigns
  (\NRunsFailed{} failed), and \NRunsSynthetic{} synthetic runs standing in for experiments not run yet.
  Every figure with synthetic data says so and draws it with hollow markers.
\end{tcolorbox}

\section{Introduction}
\label{sec:intro}

Upcoming exascale clusters will be evaluated on large-scale LLM pre-training. The work started as the
validation of two new milabench workloads at a scale they had never run at. Full access to Rorqual extends
it: the same workloads serve as probes of how the cluster, and its network architecture in particular,
behaves under large-scale training, which is what its administrators and the exascale evaluation need to
know. Rorqual is compared with Fir and Tamia, two clusters with the same GPU nodes and
four times the inter-node bandwidth per node (\cref{tab:clusters}).

\begin{table}[h]
  \caption{Clusters. All nodes have 4 H100 80\,GB GPUs connected by NVLink.}
  \label{tab:clusters}
  \centering\small
  \begin{adjustbox}{max width=\linewidth}
  \begin{tabular}{llll}
    \toprule
    Cluster & Nodes used & Inter-node network per node & Role \\
    \midrule
    Rorqual & 66 (264 GPUs), 16 leaf switches & $1\times$ 200\,Gb/s HDR InfiniBand (50\,Gb/s per GPU) & target, full-access window \\
    Fir     & leaves of 20 nodes & $4\times$ 200\,Gb/s (one per GPU) & preliminary runs \\
    Tamia   & 5 leaves of 8--16 nodes & $4\times$ 200\,Gb/s & packed vs.\ spread placement \\
    \bottomrule
  \end{tabular}
  \end{adjustbox}
\end{table}

\paragraph{Rorqual topology.} Two levels: 25 leaf switches under one spine switch. Of the 93 GPU nodes, 27
are MIG-partitioned (9 leaves) and unusable here, leaving 66 full nodes: 13 leaves with 3 GPU nodes each,
and three larger leaves with 7, 8 and 12 (ihpc218, ihpc217, ihpc316). Every leaf except ihpc316 also
connects CPU nodes, whose jobs share its uplinks.

\paragraph{Workloads.}
\begin{itemize}[leftmargin=*, itemsep=1pt]
  \item \textbf{Qwen3.8-27B} (dense, 64 layers: 48 Gated DeltaNet + 16 gated attention), TorchTitan,
        FSDP over all GPUs, $3 \times 4{,}096$ tokens per GPU per step (weak scaling), BF16 compute with
        FP32 master weights and AdamW moments, full activation checkpointing, regional compilation.
  \item \textbf{GLM-5} (mixture of experts), Megatron via slime: full model at about 208--264 GPUs,
        truncated versions at smaller scale. \placeholder{configuration once the pipeline is fixed.}
  \item \textbf{Communication benchmarks} (nccl-tests all-reduce and all-to-all) and \textbf{health
        checks} before each allocation.
\end{itemize}

\paragraph{Questions.} Each section answers one question:
(\cref{sec:hardware}) do nodes perform as expected and uniformly?
(\cref{sec:fabric}) how much does the network limit training, and does placement on leaf switches matter?
(\cref{sec:scaling}) how does throughput per GPU change from one switch to 264 GPUs?
(\cref{sec:moe}) how does an MoE model behave at scale?
(\cref{sec:software}) how does the software stack compare with expectations and published numbers?
(\cref{sec:operations}) is the system ready for production runs (startup, failures, variance, stability)?

\paragraph{Declared deviations} from the reference training recipe, valid for every result unless noted:
AdamW instead of Muon; multi-token prediction (MTP) off; random tokens instead of a corpus; FP8 not
used for Qwen (not supported on our TorchTitan pin; used for GLM-5).

\section{Methodology}
\label{sec:method}

\paragraph{Throughput.} Each run skips 30 warm-up steps (compilation, allocator growth), then measures
for 5 minutes, converted to a step count that all ranks agree on. Step times come from CUDA events on
rank~0. The official throughput is the \emph{median} of the per-step global token rates, reported per
GPU (tokens/s per GPU). Spread within a run is the 10th--90th percentile of step times.

\paragraph{MFU} $= \text{FLOPs per token} \times \text{tokens/s per GPU} / \text{peak}$, with TorchTitan's
FLOP count (6 $\times$ active parameters plus attention terms, no causal halving) and a dense BF16 peak of
989\,TFLOPS per H100. MFU is derived from a measured rate and an analytic FLOP count.

\paragraph{Placement families} (at a fixed node count):
\textbf{S}, all nodes on one switch (ihpc316, the 12-node leaf without CPU nodes);
\textbf{H}, 3 nodes on each of $N/3$ switches (3-node leaves; big leaves only beyond 13 switches);
\textbf{O}, 1 node on each of $N$ switches (up to 16);
\textbf{P}, the fewest switches (largest leaves first);
\textbf{ALL}, every healthy node.
S, H and O are compared at 3, 6, 9 and 12 nodes; at scale, H15--P15, H27--P27 and H48--P48 (variant
names give the node count). Placements are written as leaves $\times$ nodes per leaf, e.g.\ $9\times3$.

\paragraph{Expert-group placement} (GLM-5, Rorqual; an expert group is an expert-parallel group of 8
GPUs = 2 nodes). Three layouts on the same 52 nodes (208 GPUs), the largest set with an even number of
nodes on every leaf ($13\times2$ on the 3-node leaves, $12+8+6$ on the big ones): \emph{aligned}, each
group's 2 nodes on one switch; \emph{crossed}, each group's 2 nodes on two switches; \emph{mixed}, about
one group in four crossing, close to the share of a production job on every node. Megatron forms groups
from consecutive ranks, so only the node order differs: model, scale, nodes and data parallelism are the
same. A fourth run uses every healthy node in Megatron's default order (\emph{mixed, every node}): the
largest achievable job, with more nodes and some groups crossing.

Node lists are computed from the healthy nodes of the allocation by \texttt{tt\_pick\_nodes.py}.

\paragraph{Traffic isolation.} Placement comparisons are only valid between runs exposed to the same
external traffic. Three levels: \emph{alone on the cluster} (the October~30 window); \emph{reserved, one
switch} (every node in our reservation and on one leaf: other users' traffic never shares our links);
\emph{shared links} (the run crosses the spine, or its nodes are not reserved: other users' traffic may
share its links). Before the window, any run spanning switches is in the last level, so preliminary
comparisons between one switch and several mix placement with external traffic; they are not compared.

\paragraph{Health checks.} Before any run, every node of the allocation is checked (GPU count, MIG, ECC,
InfiniBand link state and rate, DCGM diagnostics, NCCL bandwidth within the node, through its NIC, in
pairs and per leaf). Failing nodes are excluded for the rest of the allocation; a node is used whole or
not at all.

\paragraph{Measured, estimated and synthetic values.}
\begin{itemize}[leftmargin=*, itemsep=1pt]
  \item \emph{Measured}: filled markers, solid lines.
  \item \emph{Estimated} (extrapolated or model-based, e.g.\ compute time scaled from a 16-layer run):
        dashed lines or hatched bars, and ``est.'' in the label.
  \item \emph{Synthetic} (placeholder for an experiment not run yet): hollow markers, and the figure says so.
\end{itemize}
Estimates never replace measurements; both are kept in the data.

\paragraph{Data.} Everything is regenerated from raw results by one command
(\texttt{scripts/analysis/analyze.py all}): milabench run data (one JSONL file per run, including failed
runs), health-check outputs, and node telemetry (\texttt{benchmate.nodemon}: GPU clocks, throttle reasons,
power, temperature and duty cycle, XID and ECC errors, NVLink, InfiniBand, CPU clocks, host memory and PCIe,
sampled once a second on every node of the job). Formats: \texttt{FORMAT.md}.

\section{Hardware baseline}
\label{sec:hardware}

\begin{question}
  \textbf{Do nodes perform as expected and uniformly?} A slow GPU or a degraded link slows every rank of
  a synchronous job, so outliers must be found before the runs, not inferred from them.
\end{question}

\Cref{fig:health-tt} shows the health checks of the preliminary allocations on Rorqual: every GPU's BF16
matmul throughput and every node's intra-node all-reduce bandwidth, grouped by leaf switch. A node is
excluded if its slowest GPU is more than 5\,\% below the allocation median, or its intra-node bandwidth
more than 10\,\% below.

\begin{figure}[h]
  \centering\fig{health_tt}
  \caption{Preliminary health checks on Rorqual (\texttt{tt\_healthcheck.sh}). Left: BF16 matmul
    ($8192^2$) per GPU; right: intra-node all-reduce bus bandwidth per node (1\,GiB). One point per
    measurement; a node checked in several allocations appears several times.}
  \label{fig:health-tt}
\end{figure}

\placeholder{spread (max/min) of matmul TFLOPS; whether intra-node bandwidth differs by leaf (two levels,
$\sim$320 and $\sim$338\,GB/s, visible above: check whether it follows the node model or the leaf).}

\Cref{fig:health-window} is the window version, from the health check that will run in the window
(srygaard/slurm-healthcheck): NVLink, NIC loopback and pairwise bandwidths for every node.

\begin{figure}[h]
  \centering\fig{health_srygaard}
  \caption{\syntheticnote{health_srygaard}Window health check on all 66 nodes, by leaf. Dotted: the
    tool's failure thresholds (absolute floors; for pairs, also 70\,\% of the median); gray: median.
    Nodes below a threshold are crossed and named.}
  \label{fig:health-window}
\end{figure}

\placeholder{number of nodes excluded and why; whether failures cluster on a leaf.}

\paragraph{Same GPU everywhere?} Cross-cluster comparisons (\cref{sec:fabric,sec:scaling}) attribute
multi-node differences to the network. That holds only if a GPU computes at the same speed on every
cluster. \Cref{fig:health-gpus} compares the per-GPU matmul throughput of the health checks.

\begin{figure}[h]
  \centering\fig[0.62\linewidth]{health_gpus}
  \caption{BF16 matmul throughput per GPU on each cluster (health checks). Black: median; dotted: 5\,\%
    below the median (exclusion threshold). \syntheticnote{health_gpus}}
  \label{fig:health-gpus}
\end{figure}

\placeholder{medians per cluster. The same holds for training: single-node runs (16 layers) on Rorqual
and Fir are within 0.6\,\% at $2\times8192$ tokens per GPU (\cref{tab:results}). The Rorqual single-node run at $3\times4096$ was noisy (p90/median 1.31) and is to be repeated: it sets Rorqual's ceiling in \cref{fig:scaling}. The window health check
(srygaard) reports each GPU's fp16 TFLOPS from the DCGM diagnostic, a coarse estimate.}

\section{Fabric and topology}
\label{sec:fabric}

\begin{question}
  \textbf{How much does the network limit training?} Does placement on leaf switches matter (same leaf
  vs.\ across leaves, nodes per switch, packed vs.\ balanced)? How does Rorqual's single HCA per node
  compare with clusters that have one per GPU?
\end{question}

\paragraph{Raw bandwidth.} \Cref{fig:comm} places each communication scope on one axis, from NVLink inside a
node to all-to-all across every leaf. On Rorqual, every all-reduce that leaves a node is capped near the
line rate of its single 200\,Gb/s HCA (25\,GB/s). The scopes are those the health checks measure:
per node, per pair, and all the nodes of a check job. The window needs two health-check runs to cover
both one leaf and the whole machine: \texttt{--job-per switch} (default) and \texttt{--job-per cluster}.
Intermediate scopes (2 or 3 leaves) come from the allocation-wide all-reduce of \texttt{tt\_healthcheck}.

\begin{figure}[h]
  \centering\fig[0.75\linewidth]{comm_scopes}
  \caption{Peak bus bandwidth by scope, over the tested message sizes (128\,MiB--1\,GiB). Gray:
    preliminary health checks (\texttt{tt\_healthcheck}); orange: uplink tests between ihpc217 and ihpc316
    (pairs alone, then 7 pairs at once; 2026-10-24); blue: window health check (srygaard). All-to-all
    and all-reduce bus bandwidths are not comparable (different normalization). \syntheticnote{comm_scopes}}
  \label{fig:comm}
\end{figure}

\placeholder{per-scope medians; the drop across leaves including ihpc316 (19.2 vs.\ 24.6\,GB/s in the
NCCL tests, $-35\,\%$ for 217+316) and whether it comes from ihpc316's uplinks.}

\paragraph{Placement at scale: balance or pack?} Up to 12 nodes, placement is analysed with scale in
\cref{sec:scaling} (\cref{fig:placement-scaling}). Beyond 12 nodes no switch holds the job, and the
allocation choice is between spreading it 3 per switch (balanced, over the 3-node leaves) and packing it
into the fewest switches (the big leaves first). \Cref{fig:placement-at-scale} plots throughput per GPU
against total throughput, so both read at once: runs with the same GPU count lie on one ray through the
origin. Along a ray, placement changes only throughput per GPU; from one ray to the next, adding nodes
lowers throughput per GPU but raises the total.

\begin{figure}[h]
  \centering\fig{placement_at_scale}
  \caption{Rorqual beyond one switch: throughput per GPU vs.\ total throughput, for 3 per switch
    (balanced) and fewest switches (packed). Gray rays: same GPU count. \syntheticnote{placement_at_scale}}
  \label{fig:placement-at-scale}
\end{figure}

\placeholder{measured so far, all preliminary: H6 = S6 within 0.2\,\% on Rorqual, repeats within 0.1\,\%; but
H6 crossed the spine outside our control while S6 stayed in one reserved switch, so the pair mixes placement
with possible traffic from other users. P12 split 6+6 across ihpc218 and ihpc316 is 4.5\,\% below S6, not
separable from 6$\to$12-node scaling.}

\paragraph{Per-GPU batch.} With FSDP, every step all-gathers the parameters and reduce-scatters the
gradients: a volume that does not depend on the batch. If communication dominates, more tokens per GPU
per step cost almost no extra time (\cref{fig:batch}).

\begin{figure}[h]
  \centering\fig{batch_effect}
  \caption{Throughput vs.\ tokens per GPU per step. On 2 nodes, all runs use BF16 AdamW moments (low-memory
    variant).}
  \label{fig:batch}
\end{figure}

\placeholder{Rorqual 2 nodes: $1\to2$ sequences of 8192 gives $+80\,\%$ (954 $\to$ 1722); Fir 3 nodes:
$+24\,\%$. Step time on Rorqual barely changes with the batch: communication-bound.}

\section{Scaling}
\label{sec:scaling}

\begin{question}
  \textbf{How does throughput per GPU change from one switch up to 264 GPUs, and where does efficiency
  drop?} The work per GPU is fixed (weak scaling), so perfect scaling is a flat line.
\end{question}

\Cref{fig:scaling} shows tokens/s per GPU against the number of GPUs, one line per cluster, with the
change from the smallest to the largest run written at the end of each line. There is no ``efficiency''
axis: the only runs without inter-node traffic are single-node runs, which use a 16-layer model, so no
full-model run can serve as a 100\,\% reference. The dashed lines are a \emph{no-network ceiling}
instead: the single-node MFU converted to full-model tokens/s. It is an estimate: it assumes MFU carries
over from 16 to 64 layers, and the 16-layer model also gathers 4$\times$ fewer parameters per step.

\begin{figure}[h]
  \centering\fig{weak_scaling}
  \caption{Weak scaling of Qwen3.8-27B ($3\times4096$ tokens per GPU, FP32 moments). One point per run
    (all placements), lines through the medians. Dashed: no inter-node traffic (est.).
    \syntheticnote{weak_scaling}}
  \label{fig:scaling}
\end{figure}

\placeholder{throughput at 12, 24, 48 GPUs (measured) and at full scale; where the curve bends.}

\paragraph{More nodes vs.\ crossing switches.} Throughput per GPU drops as a job grows (\cref{fig:scaling}).
\Cref{fig:placement-scaling} separates two causes up to 12 nodes, the largest job one Rorqual switch holds.
The one-switch line (S) is the cost of more nodes alone, with no traffic through the spine; its band is the
run-to-run spread. The spread placements on the same axis, 3 per switch (H) and 1 per switch (O), show what
crossing switches adds. At 3 nodes, 3 per switch is a single 3-node leaf, so it is not shown as spreading.
Instead, a leaf-to-leaf control (two one-switch runs on different leaves) shows how much leaves differ by
themselves. All clusters share one axis: each is relative to its own one-switch run at 3 nodes. 1 per
switch is limited by the number of leaves (16 on Rorqual, 8 GPU leaves on Fir, 5 on Tamia). Beyond 12
nodes, the two causes cannot be separated on Rorqual (no one-switch reference); the allocation choice at
that scale is in \cref{fig:placement-at-scale}.

\begin{figure}[h]
  \centering\fig{placement_scaling}
  \caption{Throughput per GPU by number of nodes for each placement (marker) and cluster (color), relative to
    the cluster's one-switch run at 3 nodes; one campaign per cluster. Band: run-to-run spread around one
    switch. Shaded (Rorqual): what crossing switches adds. Gray X: leaf-to-leaf control.
    \syntheticnote{placement_scaling}}
  \label{fig:placement-scaling}
\end{figure}

\placeholder{share of the drop at 12 nodes due to node count vs.\ crossing switches, per cluster.}

\paragraph{Where the step time goes.} \Cref{fig:decomposition} splits the step time into compute and exposed
communication. The compute part is the single-node step time at the same per-GPU batch, scaled by the
ratio of FLOPs per token between the full and the 16-layer model. It is an estimate, and so is the
communication part derived from it.

\begin{figure}[h]
  \centering\fig[0.85\linewidth]{step_decomposition}
  \caption{Step time decomposition by cluster and scale. Hatched: compute (est.); shaded: exposed
    communication (est.), the remainder of the measured step time. \syntheticnote{step_decomposition}}
  \label{fig:decomposition}
\end{figure}

\placeholder{share of exposed communication on Rorqual vs.\ Fir at equal scale; a measured split would
need gradient-accumulation runs at multi-node scale ($t_\text{step} = k\,t_\text{micro} +
t_\text{overhead}$).}

\section{Mixture-of-experts workloads}
\label{sec:moe}

\begin{question}
  \textbf{How does GLM-5 behave at scale?} What does expert traffic cost when it leaves the node or crosses
  switches, and what does FP8 bring?
\end{question}

Expert parallelism (EP) adds all-to-all exchanges every MoE layer. With EP$\le$4, experts stay inside a
node (NVLink); beyond that, all-to-all traffic crosses the single HCA of each node and, at larger EP, leaf
switches. All-to-all bandwidth across leaves is the weakest scope measured (\cref{fig:comm}).

\begin{figure}[h]
  \centering\fig{moe_scaling}
  \caption{\syntheticnote{moe_scaling}GLM-5 MFU (on active parameters) vs.\ GPUs, by expert-parallel size, BF16 and
    FP8. Left of the dotted line: truncated models.}
  \label{fig:moe}
\end{figure}

\paragraph{Expert-group placement.} Keeping every EP group inside a switch (aligned) is the cheapest layout,
but it fits on only 52 of the 66 nodes (an even number per leaf). The question is whether using every node
pays off: more GPUs, but some groups crossing switches. \Cref{fig:expert-placement} (left) puts every run
on one plane, with the same scale on both axes: more GPUs than aligned (x) against more total throughput
than aligned, summed over all GPUs (y). The diagonal is throughput proportional to GPUs; below zero, a run
delivers less in total than aligned. On the same 52 nodes (x = 0), only the node order changes, which gives
the cost of a minority of crossing groups (mixed) and of the worst case (crossed). The run on every node is
worth it only if it lands above zero; with $N$ nodes, that requires a per-GPU cost below $1 - 52/N$,
i.e.\ 19\,\% at 64 nodes. The right panel shows where the extra GPUs go: the gain the extra nodes would
bring if throughput were proportional, the cost of crossing groups (the per-GPU cost of mixed vs.\ aligned
on the same nodes, applied to every node: an estimate) and the rest (larger scale, more simultaneous
crossings: the remainder, an estimate too).

\begin{figure}[h]
  \centering\fig{expert_placement}
  \caption{GLM-5 expert-group placement on Rorqual (EP8). Left: more GPUs vs.\ more total throughput than
    the largest aligned run (52 nodes), equal scales; dotted: distance below proportional scaling. Right:
    total throughput of the run on every node vs.\ aligned (= 100), step by step; solid: measured or
    arithmetic, hatched: estimated. \syntheticnote{expert_placement}}
  \label{fig:expert-placement}
\end{figure}

Every rank waits for the slowest group at each all-to-all, so a few crossing groups may cost almost as
much as all of them. \Cref{fig:expert-groups} shows the all-to-all time of each group in the mixed runs.

\begin{figure}[h]
  \centering\fig[0.6\linewidth]{expert_groups}
  \caption{All-to-all time per EP group in the mixed runs, inside one switch vs.\ across switches.
    \syntheticnote{expert_groups}}
  \label{fig:expert-groups}
\end{figure}

\placeholder{cost of mixed and crossed; total-throughput gain of every node over aligned, and whether it
clears break-even.}

\placeholder{GLM-5 configuration (layers, experts, top-$k$, EP and routing modes), output format from the
Megatron pipeline (see \texttt{FORMAT.md}), and results. Per-rank expert load, if logged, would show
routing imbalance.}

\section{Software stack}
\label{sec:software}

\begin{question}
  \textbf{Is throughput where it should be?} Compared with expectations (single-node MFU, published
  numbers for similar models and hardware), with the declared deviations accounted for.
\end{question}

\placeholder{table: our MFU vs.\ published MFU for dense models of similar size on H100 (TorchTitan,
Megatron, vendor reports), with each source's batch, sequence length, precision and parallelism. Effect of
each declared deviation: Muon vs.\ AdamW (optimizer step cost), MTP (extra FLOPs, an analytic correction
shown as an estimate), FP8 for GLM-5.}

\placeholder{software versions: TorchTitan 51b37c2, PyTorch 2.15.0.dev20261005+cu130, NCCL version,
compile mode and its startup cost.}

\section{Operational readiness}
\label{sec:operations}

\begin{question}
  \textbf{Is the system ready for production training?} How long until training starts, how often do runs
  fail, how much do repeated runs vary, and is throughput stable over an hour or more?
\end{question}

\paragraph{Startup.} \Cref{fig:startup} separates process start to an initialized model (imports, process
groups, model construction and sharding) from the first step (compilation, or loading the compile cache).

\begin{figure}[h]
  \centering\fig{startup}
  \caption{Startup phases by number of nodes. \syntheticnote{startup}}
  \label{fig:startup}
\end{figure}

\placeholder{outliers are explained: 297\,s init with the virtual environment on Lustre instead of node-local
disk; first steps of 85--117\,s with a cold compile cache (new shapes).}

\paragraph{Stability.} \Cref{fig:stability} follows a long run step by step, with the GPU clocks from DCGM on
the same time axis, so that a throughput drop can be traced to a GPU.

\begin{figure}[h]
  \centering\fig{stability}
  \caption{\syntheticnote{stability}Throughput per step over one hour on every healthy node (top) and SM clocks of
    all GPUs (bottom): the gray band is the range of the other GPUs, the colored line the GPU with the
    lowest clock.}
  \label{fig:stability}
\end{figure}

\paragraph{Stragglers.} With FSDP, all ranks wait for the slowest one, so per-rank step times are identical
(measured: 24 ranks within 0.2\,ms) and cannot identify a slow GPU. The slow GPU is instead the one that
\emph{waits least}: its tensor cores are busy for a larger share of each step (\cref{fig:stragglers}).

\begin{figure}[h]
  \centering\fig{stragglers}
  \caption{\syntheticnote{stragglers}Tensor-pipe activity per GPU (DCGM), the highest 10-minute median of each GPU,
    by leaf. The outlier is named.}
  \label{fig:stragglers}
\end{figure}

\placeholder{run-to-run variance (repetitions in \cref{fig:placement-scaling,fig:placement-at-scale}; so far $\le$0.1\,\% on Rorqual within an
allocation), failures and their causes (\cref{tab:ledger}), excluded nodes.}

\section{Conclusions and recommendations}
\label{sec:conclusion}

\placeholder{answers to the six questions, in order; recommendations for the Rorqual administrators
(network, node placement, health checks) and for the evaluation of the exascale clusters (what to measure,
which workloads, which placement experiments).}


\appendix
\section{Window experiment planning}
\label{app:planning}

The full-access window (expected 2026-10-29 17:00 to 10-30 17:00) is planned so that the most important
results come first and the plan degrades gracefully if the window is cut short. The detailed checklist,
with the run-by-run queue and owners, is in \texttt{scripts/analysis/EXPERIMENTS.md}.

\paragraph{Prioritization.} The campaign started as a validation of two new milabench workloads (Qwen3.8-27B
on TorchTitan and GLM-5 on Megatron) at a scale they had never run at. Full access to Rorqual turned it into
an opportunity to also use these workloads as probes of the questions that matter to the cluster's
administrators. The ranking keeps both purposes in view, with four criteria: whether a result shows that
each workload runs and measures correctly at full scale (the original purpose), how directly it answers the
cluster questions (does Rorqual's network limit large-scale training, and what follows for allocation and
for the exascale evaluation), whether it informs a decision, and whether only the window can produce its
data. The first criterion does not change the order: the tier-1 runs already include each workload on every
healthy node. Two rules refine the ranking:
\begin{itemize}[leftmargin=*, itemsep=1pt]
  \item \textbf{Tier 0 before anything else}: the health check of every node is a prerequisite, not a
        result. It sets the healthy set, the exclusions and the node lists of every later run, and it is
        the first snapshot of the state of the cluster.
  \item \textbf{GLM-5 first when near-equal}: GLM-5 runs are more complex and more likely to fail; running
        them first leaves more time to fix them before the window ends.
\end{itemize}

\begin{table}[h]
  \caption{Figure priority. Tier 1 alone answers the main questions; tier 2 explains them; tier 3 is
    context, or already covered by preliminary data.}
  \label{tab:priority}
  \centering\small
  \begin{tabular}{cll}
    \toprule
    Tier & Figure & Role \\
    \midrule
    0 & Window health check (data) & sanity check, state of the cluster, node lists \\
    \midrule
    1 & \Cref{fig:expert-placement} GLM-5 expert placement & using every node vs.\ the largest aligned subset \\
    1 & \Cref{fig:scaling} Weak scaling & throughput per GPU up to the full machine \\
    1 & \Cref{fig:placement-at-scale} At scale: balance or pack & allocation of large jobs \\
    1 & \Cref{fig:placement-scaling} More nodes vs.\ crossing switches & where the scaling loss comes from \\
    \midrule
    2 & \Cref{fig:comm} Bandwidth by scope & physical cause (one HCA per node, spine) \\
    2 & \Cref{fig:health-window} Window health check & validity of the rest \\
    2 & \Cref{fig:moe} GLM-5 scaling, EP size, FP8 & MoE at scale \\
    2 & \Cref{fig:expert-groups} All-to-all per EP group & mechanism behind \cref{fig:expert-placement} \\
    2 & \Cref{fig:decomposition} Step decomposition & communication-bound step \\
    \midrule
    3 & \Cref{fig:batch} Per-GPU batch & answered by preliminary data \\
    3 & \Cref{fig:startup} Startup & fed by every run \\
    3 & \Cref{fig:stragglers} Stragglers & node telemetry from regular runs \\
    3 & \Cref{fig:stability} Stability & conditional on window length \\
    3 & \Cref{fig:health-gpus} Per-GPU compute by cluster & premise shown by preliminary data \\
    3 & \Cref{fig:health-tt} Preliminary health checks & background \\
    \bottomrule
  \end{tabular}
\end{table}

\paragraph{Roles.} Three people run the window: one for health checks and monitoring (runs the health
checks, watches GPU telemetry and logs during every run, regenerates the quick figures after each block and
compares them with expectations, flags anything suspicious), one for the Qwen workload and one for the GLM-5
workload.

\paragraph{Rules.}
\begin{itemize}[leftmargin=*, itemsep=1pt]
  \item \textbf{Small blocks, alternating workloads}: each block is one comparison (10--36 minutes) and
        launches from a prepared command; GLM-5 and Qwen blocks alternate, so one team prepares while the
        other runs.
  \item \textbf{Skip on failure}: a failed or suspicious block is logged and the next block of the
        \emph{other} workload starts; the failed block returns to the queue once fixed.
  \item \textbf{First pass, then repeats}: pass 1 runs every tier-0 and tier-1 comparison once, so every key
        figure has data early; pass 2 repeats it; then the remaining points and the optional runs.
  \item \textbf{Parallel runs}: 1-node runs (smoke tests, re-checks) use idle nodes at any time; a run that
        stays on one switch runs next to any run that does not use that switch; runs that cross the spine run
        one at a time. The assumption that a switch does not slow traffic between its own nodes is checked in
        the window (the same one-switch run with and without a concurrent run elsewhere).
  \item \textbf{Exclusions}: a node flagged by the health check is excluded and the node lists are
        recomputed; a comparison interrupted by an exclusion is rerun on the new node set.
\end{itemize}

\begin{table}[h]
  \caption{Window phases (Qwen runs measured at $\sim$12 minutes; GLM-5 runs assumed at $\sim$15 minutes until
    measured), with parallel runs. About 11.5 hours of execution, leaving $\sim$12.5 hours of the window for
    fixes and failures.}
  \label{tab:phases}
  \centering\small
  \begin{tabular}{lp{0.72\linewidth}r}
    \toprule
    Phase & Content & Time \\
    \midrule
    Pass 1 & health check; GLM-5 aligned, crossed, mixed, every node; Qwen every node, H27/P27, S3/S12/O12, H12,
             leaf-to-leaf controls; whole-machine health checks & 2.5 h \\
    Pass 2 & second repetition of pass 1 (one-switch check run alone); health re-check & 2.4 h \\
    P2 & GLM-5 FP8; Qwen H15/P15, H48/P48, 6- and 9-node placements, O16, O3; their repeats & 4.5 h \\
    P3 & GLM-5 EP16; bridge to past data with gradient accumulation; 1-hour stability run; final health check & 2.1 h \\
    \bottomrule
  \end{tabular}
\end{table}

\paragraph{Coverage.} \Cref{tab:coverage} states, for each figure, what the report would contain if the
window stopped after each phase.

\begin{table}[h]
  \caption{Report coverage by phase (cumulative). \emph{none}: no window data; \emph{usable}: supports a
    statement, with the limit given; \emph{complete}: as designed.}
  \label{tab:coverage}
  \centering\footnotesize
  \setlength{\tabcolsep}{3pt}
  \begin{tabular}{c>{\raggedright}p{3.1cm}>{\raggedright}p{2.35cm}>{\raggedright}p{2.35cm}>{\raggedright}p{2.2cm}>{\raggedright\arraybackslash}p{2.0cm}}
    \toprule
    Tier & Figure & After pass 1 & After pass 2 & After P2 & After P3 \\
    \midrule
    0 & Health check (data) & complete & complete & complete & complete \\
    1 & \ref{fig:expert-placement} Expert placement & usable: 1 run each, no significance & \textbf{complete} & complete & complete \\
    1 & \ref{fig:scaling} Weak scaling & usable: 12--264 GPUs, 5 scales & usable, repeats & \textbf{complete} & complete \\
    1 & \ref{fig:placement-at-scale} Balance or pack & usable: 27 nodes, 1 run & usable: with run-to-run spread & \textbf{complete}: 15, 27, 48 & complete \\
    1 & \ref{fig:placement-scaling} Nodes vs.\ crossing & usable: 3 and 12 nodes, controls & usable, repeats & \textbf{complete}: 6, 9, 16 & complete \\
    2 & \ref{fig:comm} Bandwidth by scope & \textbf{complete} & complete & complete & complete \\
    2 & \ref{fig:health-window} Window health check & \textbf{complete} & complete & complete & complete \\
    2 & \ref{fig:moe} GLM-5 scaling & usable: EP8 BF16 & usable, repeats & usable: + FP8 & \textbf{complete}: + EP16 \\
    2 & \ref{fig:expert-groups} All-to-all per group & complete if timers logged & complete & complete & complete \\
    2 & \ref{fig:decomposition} Step decomposition & usable (estimated) & usable & usable & \textbf{complete} if measured \\
    3 & \ref{fig:batch} Per-GPU batch & usable (preliminary) & same & same & \textbf{complete} \\
    3 & \ref{fig:startup} Startup & usable & usable & \textbf{complete} & complete \\
    3 & \ref{fig:stragglers} Stragglers & usable (5-min windows) & usable & usable & \textbf{complete} \\
    3 & \ref{fig:stability} Stability & none & none & none & \textbf{complete} \\
    3 & \ref{fig:health-gpus} Per-GPU compute & complete if tool reports it & same & same & same \\
    3 & \ref{fig:health-tt} Preliminary health checks & complete & complete & complete & complete \\
    \bottomrule
  \end{tabular}
\end{table}

\paragraph{What is lost if the window stops early.} Before pass 2 ends, the placement results
(\cref{fig:expert-placement,fig:placement-at-scale}) have no significance test: differences of 0.2--1\,\% cannot
be told from run-to-run noise with one run each. Before P2 ends, the tier-1 curves stay sparse and GLM-5 has
no FP8 result. Before P3, there is no stability figure, the step decomposition stays estimated, GLM-5 has no
EP-size comparison at full scale, and the batch figure has no window point. Independently of the window,
\cref{fig:expert-groups} depends on a tool change made before it (per-group all-to-all timers in Megatron).

\section{Results}
\label{app:results}

\begin{table}[h]
  \caption{Completed runs (measured only). L: layers; tok/GPU: sequences $\times$ length ($\times$
    accumulation); leaves: nodes per leaf switch; mom.: AdamW moment dtype; p90/p10: step-time spread;
    GiB: peak allocated / reserved; init and 1st: startup phases (s).}
  \label{tab:results}
  \centering
  \begin{adjustbox}{max width=\linewidth}
    \input{generated/tables/results}
  \end{adjustbox}
\end{table}

\begin{table}[h]
  \caption{Run ledger: every attempt since the current measurement path (2026-10-21), including failures.}
  \label{tab:ledger}
  \centering
  \begin{adjustbox}{max width=\linewidth}
    \input{generated/tables/ledger}
  \end{adjustbox}
\end{table}


\end{document}
